% from ega4-11-fr.tex
\oldpage[IV-3]{154}

Pour appliquer ce lemme à la situation de (11.5.5) on notera en outre que si
l'homomorphisme $A\to A'$ est injectif, il en est ainsi \emph{a fortiori} de
$A\to A'_{\lambda}$, pour tout $\lambda$.

\begin{lemma}[11.5.5.2]
\label{IV.11.5.5.2-fr}
Soient $A$ un anneau, $A'$ une $A$-algèbre, $(A'_{\lambda})$ un système inductif
de $A$-algèbres tel que $A'=\varinjlim A'_{\lambda}$~; on pose
$Y=\operatorname{Spec}(A)$, $Y'=\operatorname{Spec}(A')$,
$Y'_{\lambda}=\operatorname{Spec}(A'_{\lambda})$. Soient $f:X\to Y$ un
morphisme de présentation finie, $\mathcal F$ un $\mathcal O_X$-Module
quasi-cohérent de présentation finie~; on pose $X'=X\times_Y Y'$,
$f'=f_{(Y')}$, $\mathcal F'=\mathcal F\otimes_Y Y'$,
$X'_{\lambda}=X\times_Y Y'_{\lambda}$, $f'_{\lambda}=f_{(Y'_{\lambda})}$,
$\mathcal F'_{\lambda}=\mathcal F\otimes_Y Y'_{\lambda}$. Soit $x$ un point de
$X$, tel que $\mathcal F'$ soit $f'$-plat en tous les points $x'\in X'$
au-dessus de $x$~; alors il existe $\lambda$ tel que $\mathcal F'_{\lambda}$
soit $f'_{\lambda}$-plat en tous les points de $X'_{\lambda}$ au-dessus de $x$.
\end{lemma}

\begin{proof}
Soit $U'$ l'ensemble des $x'\in X'$ tels que $\mathcal F'$ soit $f'$-plat au
point $x'$~; on sait (11.3.1) que $U'$ est ouvert dans $X'$ puisque $f'$ est de
présentation finie (1.6.2)~; de même l'ensemble $U'_{\lambda}$ des points de
$X'_{\lambda}$ où $\mathcal F'_{\lambda}$ est $f'_{\lambda}$-plat est ouvert
dans $X'_{\lambda}$, et on sait en outre (11.2.6) que $U'$ est réunion des
$v_{\lambda}^{-1}(U'_{\lambda})$, où
$v_{\lambda}:X'\to X'_{\lambda}$ est la projection canonique. Considérons le
schéma $T=\operatorname{Spec}(k(x))$~; posons $T'=T\times_Y Y'$,
$T'_{\lambda}=T\times_Y Y'_{\lambda}$, et soient
$w_{\lambda}:T'\to T'_{\lambda}$, $g':T'\to X'$,
$g'_{\lambda}:T'_{\lambda}\to X'_{\lambda}$ les projections canoniques. Posons
$V'_{\lambda}=g_{\lambda}'^{-1}(U'_{\lambda})$,
$V'=g'^{-1}(U')=\bigcup_{\lambda}w_{\lambda}^{-1}(V'_{\lambda})$. Par
hypothèse on a (compte tenu de (\textbf I, 3.6.1)) $V'=T'$~; comme $T'$ est
quasi-compact, il existe $\lambda$ tel que
$w_{\lambda}^{-1}(V'_{\lambda})=T'$. On déduit alors de (8.3.3) appliqué aux
parties fermées quasi-compactes $T'_{\lambda}-V'_{\lambda}$ de
$T'_{\lambda}$, qu'il existe $\mu\geq\lambda$ tel que
$T'_{\mu}=V'_{\mu}$~; cela signifie que $\mathcal F'_{\mu}$ est
$f'_{\mu}$-plat en tous les points de $X'_{\mu}$ dont la projection dans $X$
est $x$. C.Q.F.D.
\end{proof}

\subsection{Descente de la platitude par des morphismes quelconques : cas d'un
préschéma de base unibranche}
\label{subsection:IV.11.6-fr}

\begin{theorem}[11.6.1]
\label{IV.11.6.1-fr}
Soient $A$ un anneau local intègre géométriquement unibranche
$(\mathbf 0,23.2.1)$, $Y=\operatorname{Spec}(A)$, $f:X\to Y$ un morphisme
localement de présentation finie, $\mathcal F$ un $\mathcal O_X$-Module
quasi-cohérent de présentation finie. Soient $A'$ un anneau local,
$u:A\to A'$ un homomorphisme local injectif~; on pose
$Y'=\operatorname{Spec}(A')$, $X'=X\times_Y Y'$, $f'=f_{(Y')}$,
$\mathcal F'=\mathcal F\otimes_Y Y'$. Soient $x$ un point de $X$ dont la
projection $f(x)=y$ est le point fermé de $Y$, $x'$ un point de $X'$ dont les
projections dans $X$ et $Y'$ sont respectivement $x$ et le point fermé $y'$ de
$Y'$. Alors, si $\mathcal F'$ est $f'$-plat au point $x'$, $\mathcal F$ est
$f$-plat au point $x$.
\end{theorem}

\begin{proof}
On peut se borner au cas où $f$ est de présentation finie, la question étant
locale sur $X$. Nous procéderons en plusieurs étapes.

\medskip
\noindent\textbf{I)} \emph{Réduction au cas où $A$ et $A'$ sont des anneaux
locaux intégralement clos.}

Comme $u$ est injectif et $A$ intègre, il existe un idéal premier
$\mathfrak p'$ de $A'$ tel que $u^{-1}(\mathfrak p')=0$
$(\mathbf 0_{\mathrm I},1.5.8)$~; l'homomorphisme composé
$A\xrightarrow{u}A'\to A''=A'/\mathfrak p'$ est donc injectif et local, et si
$Y''=\operatorname{Spec}(A'')$, $X''=X'\otimes_{A'}A''=X\times_Y Y''$,
$\mathcal F''=\mathcal F'\otimes_{A'}A''$, $\mathcal F''$ est $Y''$-plat aux
points de $X''$ au-dessus de $x'$ (2.1.4)~; remplaçant au besoin $A'$ par $A''$
et tenant compte de (\textbf I, 3.4.7), on peut donc supposer d'abord que $A'$
est intègre. Si $K'$ est le corps des fractions de $A'$, il y a alors un anneau
de valuation $B'$ dans $K'$ qui domine $A'$~; l'homomorphisme composé
$A\xrightarrow{u}A'\to B'$ étant injectif et local, le même raisonnement que
% IV3_B23_P154_END

